\section{Proofs and interpolants inside admissible sets}
\label{sec:p03-admissible-proof-closure}

The model-existence and interpolation arguments of
\cref{thm:p02-w02c-p02-003-model-existence,thm:p02-w02c-p02-003-interpolation}
produce countable infinitary objects. For an admissible fragment the
stronger requirement is that the proof code and interpolant belong to
the specified admissible set. This section establishes that requirement
using Collection. It distinguishes an internally available family from
an arbitrary countable family in the surrounding universe.
The resulting admissible interpolation theorem is the ingredient used
in Barwise's abstract-logic argument
\cite[Part~II, \S4, p.~261]{Barwise1974Axioms}.

\begin{hbconvention}{def:p03-adm-setting}{The admissible setting}
Let \(A\) be a countable transitive admissible set in the sense of
\cref{def:p03-w03-p03-039-nadel-kp}, with \(\omega\in A\).
Countability is external. All schemas and definability assertions below
are relative to the actual membership structure on \(A\).
The terms \(\Delta_0\), \(\Sigma_1\), and \(\Sigma\) have the meanings
of \cref{def:p03-w03-p03-039-nadel-levy}. A \(\Delta_1\) relation is one
having complementary \(\Sigma_1\) definitions over this structure.
Parameters from \(A\) are permitted.

Formula codes are the usual tagged set codes: atomic formulas have
finite term records, a quantified formula records its variable and
matrix, and an infinitary conjunction or disjunction records its set of
immediate constituent formulas. Each formula has only finitely many free
variables. Formula codes in \(A\) constitute \(L_A\).
We first use a relational signature \(K\in A\) and a countably infinite
set \(C\in A\) of fresh constant names, disjoint from \(K\). Every
formula used in a finite condition has only finitely many names from
\(C\), including occurrences inside its constituent formulas.
The original sentences for \(K\) have no such names. Names and variables
can be coded by tagged natural numbers, with disjoint tags for the two
stocks.
\end{hbconvention}

\subsection{Collection and syntactic closure}

\begin{hblemma}{thm:p03-adm-collection}{proved=annex:p3.adm.1}{Collecting definable witnesses}
In \(\mathrm{KP}\):
\begin{enumerate}
\item every \(\Sigma\)-formula is equivalent to a \(\Sigma_1\)-formula;
\item Collection holds for \(\Sigma\)-formulas;
\item Separation holds for relations with complementary
      \(\Sigma_1\) definitions;
\item a total functional \(\Sigma\)-relation on a set has a set as its
      graph and as its range.
\end{enumerate}
In particular, if \(a\in A\) and for each \(x\in a\) there is a witness
\(y\in A\) to a \(\Sigma\)-condition, a set in \(A\) collects witnesses
for all \(x\). No selection of exactly one witness per \(x\) is asserted
unless the relation is functional.
\end{hblemma}

\begin{hblemma}{thm:p03-adm-recursion}{proved=annex:p3.adm.2}{Recursion on a set relation}
Let \(I,R\in A\), with \(R\subseteq I\times I\) externally well founded,
and put \(I_x=\{y\in I:yRx\}\). Suppose a \(\Sigma\)-formula
\(\gamma(x,g,z)\) defines, over \(A\), exactly one value \(z\in A\)
for every \(x\in I\) and every function \(g\in A\) with domain \(I_x\).
Then a unique function \(f\in A\), with domain \(I\), satisfies
\[
 A\models\gamma\bigl(x,f\restriction I_x,f(x)\bigr)
 \qquad(x\in I).
\]
The graph of the recursively assigned value is \(\Sigma\)-definable,
uniformly from \(I,R,\gamma\), on inputs satisfying these hypotheses.
No uniform test for well-foundedness of an arbitrary relation is asserted.
\end{hblemma}

\begin{hbtheorem}{thm:p03-adm-membership-recursion}{proved=annex:p3.adm.8}{Membership recursion in KPU}
Let \(\gamma(\bar x,y,g,w,\bar r)\) be a \(\Sigma\)-formula defining a
total functional operation \(G(\bar x,y,g)\). In \(\mathrm{KPU}\),
transitive closures exist and there is a unique definable operation
\(F\) satisfying
\[
 F(\bar x,y)=G\bigl(\bar x,y,
                  F_{\bar x}\restriction\mathrm{TC}(y)\bigr).
\]
Its graph is the \(\Sigma\)-formula \(\Phi_\gamma\) displayed in
\cref{thm:p02-w02c-p02-003-hf-sigma-recursion}, with the same bounded
pair, function, closure and restriction formulas. This is a theorem
schema provable in \(\mathrm{KPU}\) from the totality and functionality
of \(\gamma\); it does not require Infinity, countability, or external
well-foundedness of an arbitrary model of the axioms. It extends the
previously proved hereditarily finite instance to the stated axiom system.
\end{hbtheorem}

\begin{hbtheorem}{thm:p03-adm-truncation}{proved=annex:p3.adm.9}{The well-founded part}
Let \(\mathfrak B\) be a set-sized model of \(\mathrm{KPU}\), possibly
externally ill founded. Its well-founded part, including all its
urelements, collapses to an admissible set over those urelements.
If \(\mathfrak B\) contains a set of all its urelements, that set also
belongs to its well-founded part. If its internal \(\omega\) is standard,
the collapse contains the actual \(\omega\) as an element.
These conclusions do not assert that an arbitrary nonstandard internal
\(\omega\) survives the collapse.
\end{hbtheorem}

\begin{hblemma}{thm:p03-adm-syntax}{proved=annex:p3.adm.3}{Closure of formula codes}
The following operations preserve membership in \(A\) and have
\(\Delta_1\) graphs on their stated domains: forming finite lists,
finite sets, pairs and products; transitive closure; obtaining the
subformulas, free variables and occurring symbols of a formula;
negation normal form; capture-free substitution of a constant for a free
variable; injective renaming of variables by a map in \(A\); and the relational-image
translations of \cref{def:p02-w02c-p02-003-relational-image}.
If \(\Phi\in A\) is a set of formula codes with a common finite bound
on their free variables, its conjunction and disjunction are also in
\(A\). This last assertion requires \(\Phi\in A\), not merely
\(\Phi\subseteq A\) and external countability.
\end{hblemma}

\subsection{Ranked refutations}

\begin{hbdefinition}{def:p03-adm-refutation}{The refutation rules}
A condition \(\Gamma\) is a finite set of closed negation-normal-form
formulas in \(A\), for \(K\cup C\), with finitely many occurring names
from \(C\). A refutation of \(\Gamma\) establishes that its members
cannot all hold in a nonempty structure. Its initial conditions contain
an atomic sentence and its negation. Its rules are the following,
read from the displayed premises to the conclusion \(\Gamma\).
\begin{enumerate}
\item From a refutation of \(\Gamma\cup\{c=c\}\), infer a refutation
      of \(\Gamma\), for \(c\in C\).
\item If \(c=d\) and \(\alpha(c/v)\) belong to \(\Gamma\), where
      \(\alpha\) is atomic with at most \(v\) free, use the premise
      \(\Gamma\cup\{\alpha(d/v)\}\).
\item If \(\bigwedge\Phi\in\Gamma\), use the premise
      \(\Gamma\cup\{\chi\}\) for any chosen \(\chi\in\Phi\).
\item If \(\forall v\,\chi\in\Gamma\), use the premise
      \(\Gamma\cup\{\chi(c/v)\}\) for any \(c\in C\).
\item If \(\bigvee\Phi\in\Gamma\), use premises
      \(\Gamma\cup\{\chi\}\) for every \(\chi\in\Phi\).
      When \(\Phi\) is empty this is a rule with no premises.
\item If \(\exists v\,\chi\in\Gamma\), use the premise
      \(\Gamma\cup\{\chi(c/v)\}\), where \(c\in C\) occurs in no
      member of \(\Gamma\).
\end{enumerate}
A proof code is a set of nodes with a distinguished root, a condition
and rule record at each node, a child relation, and an ordinal-valued
rank function strictly decreasing from a node to each child. For the
fifth rule the child conditions cover every required premise; repeated
copies of a premise are allowed. For the other rules a designated child
has the indicated premise. All nodes satisfy one of the stated local
conditions. Unused nodes can be discarded but do not affect validity.
An \emph{\(A\)-refutation} is such a code belonging to \(A\).
\end{hbdefinition}

\begin{hblemma}{thm:p03-adm-proof-codes}{proved=annex:p3.adm.4}{Checking and joining proofs}
The assertion that a condition has an \(A\)-refutation is
\(\Sigma_1\) over \(A\). The refutation rules are sound. If
\(\Phi\in A\), \(\bigvee\Phi\in\Gamma\), and every
\(\Gamma\cup\{\chi\}\), \(\chi\in\Phi\), has an \(A\)-refutation,
then \(\Gamma\) has an \(A\)-refutation. The analogous closure holds
for each one-premise rule.
\end{hblemma}

\begin{hbtheorem}{thm:p03-adm-completeness}{proved=annex:p3.adm.5}{Completeness}
A condition has an \(A\)-refutation if and only if it is unsatisfiable.
Consequently validity of sentences of \(L_A\) is \(\Sigma_1\)-definable
over \(A\). The latter assertion also holds for signatures with
function and constant symbols, using their relational images.
This is an internal proof-existence statement, not a claim that the set
of valid sentences is computably enumerable without access to \(A\).
\end{hbtheorem}

\subsection{Interpolation in the same fragment}

\begin{hbtheorem}{thm:p03-adm-interpolation}{proved=annex:p3.adm.6}{Admissible interpolation}
Let \(\varphi,\psi\in L_A\) be sentences and suppose
\(\varphi\to\psi\) is valid. There is an interpolating sentence
\(\theta\in L_A\) such that \(\varphi\to\theta\) and
\(\theta\to\psi\) are valid and every non-logical symbol of \(\theta\)
occurs both in \(\varphi\) and in \(\psi\).
In particular the interpolant belongs to the given admissible set,
not just to the larger language \(L_{\omega_1\omega}\).
\end{hbtheorem}

\begin{hbcorollary}{thm:p03-adm-beth}{proved=annex:p3.adm.7}{Admissible explicit definitions}
Suppose \(\varphi(P)\in L_A\) implicitly defines a new finite-arity
relation \(P\): every reduct has at most one expansion satisfying
\(\varphi(P)\). Then an explicit defining formula for \(P\) can be
chosen in \(L_A\), using only the other symbols occurring in
\(\varphi\). This is the ordinary, possibly partial-existence, Beth
property. It is not identified here with the property for a projectively
defined relation having exactly one expansion on every reduct.
\end{hbcorollary}
